\title{Επισκόπηση Λειτουργίας Συστήματος Μεταφοράς Ηλεκτρικής Ενέργειας}

\section{Θεωρία}
\label{kef-01-theoria}

Το κεφάλαιο αυτό πραγματεύεται τους τρόπους με τους οποίους ένα σύστημα ηλεκτρικής ενέργειας
διατηρείται σε ασφαλή λειτουργία. Η λειτουργία του συστήματος μεταφοράς οργανώνεται ως
κλειστός βρόχος τεσσάρων σταδίων: συλλογή μετρήσεων, εκτίμηση της κατάστασης του δικτύου,
αξιολόγηση της επιχειρησιακής ασφάλειας και εφαρμογή ενεργειών ελέγχου. Οι ενότητες που
ακολουθούν εξετάζουν διαδοχικά τα στάδια αυτά, καθώς και τους μηχανισμούς αγοράς και
ρύθμισης συχνότητας μέσω των οποίων υλοποιούνται οι αποφάσεις λειτουργίας.

Η ορολογία ακολουθεί την επίσημη ελληνική απόδοση των Κανονισμών (ΕΕ) 2017/1485 και
2017/2195, η οποία χρησιμοποιείται και από τον ΑΔΜΗΕ. Όπου η ελληνική πανεπιστημιακή
βιβλιογραφία χρησιμοποιεί διαφορετικό όρο, αυτός παρατίθεται σε παρένθεση κατά την πρώτη
εμφάνιση, μαζί με τον αγγλικό όρο και το καθιερωμένο ακρωνύμιο.

\subsection{Ρόλος και δομή του Συστήματος Μεταφοράς}

Το Σύστημα Μεταφοράς Ηλεκτρικής Ενέργειας (ΣΜΗΕ) είναι το δίκτυο υπερυψηλής και υψηλής
τάσης που διασυνδέει τους σταθμούς παραγωγής με τα κέντρα κατανάλωσης και με τα γειτονικά
συστήματα. Το Ελληνικό Σύστημα Μεταφοράς Ηλεκτρικής Ενέργειας (ΕΣΜΗΕ) λειτουργεί στα επίπεδα
τάσης 400 kV, 150 kV και 66 kV. Διαχειριστής του είναι ο Ανεξάρτητος Διαχειριστής Μεταφοράς
Ηλεκτρικής Ενέργειας (ΑΔΜΗΕ), μέλος του Ευρωπαϊκού Δικτύου Διαχειριστών Συστημάτων Μεταφοράς
Ηλεκτρικής Ενέργειας (ΕΔΔΣΜ ηλεκτρικής ενέργειας, ENTSO-E). Το σύστημα ανήκει στη
συγχρονισμένη περιοχή της Ηπειρωτικής Ευρώπης (Continental Europe, CE).

Η μεταφορά σε υψηλή τάση περιορίζει τις απώλειες. Για δεδομένη μεταφερόμενη ισχύ
$P = \sqrt{3}\,V I \cos\varphi$, το ρεύμα είναι αντιστρόφως ανάλογο της τάσης, ενώ οι
απώλειες Joule στους αγωγούς είναι ανάλογες του $I^2$· ο διπλασιασμός της τάσης
υποτετραπλασιάζει επομένως τις απώλειες. Για τον λόγο αυτόν το σύστημα οργανώνεται
ιεραρχικά, με υπερυψηλή τάση για τη μεταφορά σε μεγάλες αποστάσεις και διαδοχικούς
υποβιβασμούς τάσης προς τα δίκτυα διανομής.

\subsubsection*{Απαιτήσεις λειτουργίας}

Ο Διαχειριστής Συστήματος Μεταφοράς (ΔΣΜ· Transmission System Operator, TSO) οφείλει να
ικανοποιεί ταυτόχρονα τις ακόλουθες απαιτήσεις:

\begin{itemize}
\item \textbf{Ισοζύγιο ισχύος.} Η ηλεκτρική ενέργεια δεν αποθηκεύεται σε σημαντικές
      ποσότητες εντός του συστήματος, οπότε η παραγωγή πρέπει να ισούται συνεχώς με το
      άθροισμα της ζήτησης και των απωλειών. Κάθε ανισορροπία καλύπτεται αρχικά από την
      κινητική ενέργεια των στρεφόμενων μαζών και εκδηλώνεται ως μεταβολή της συχνότητας:
      έλλειμμα παραγωγής προκαλεί επιβράδυνση των σύγχρονων μηχανών και πτώση της
      συχνότητας.
\item \textbf{Επιχειρησιακή ασφάλεια.} Το σύστημα πρέπει να παραμένει εντός των ορίων
      επιχειρησιακής ασφάλειας (operational security limits) — θερμικά όρια του
      εξοπλισμού, όρια τάσης, όρια ευστάθειας — τόσο στην τρέχουσα κατάσταση όσο και μετά
      από οποιοδήποτε μεμονωμένο απρόβλεπτο συμβάν (contingency), όπως η απώλεια γραμμής ή
      μονάδας. Η απαίτηση αυτή διατυπώνεται ως κριτήριο N-1, το οποίο ορίζεται στην επόμενη
      υποενότητα.
\item \textbf{Ποιότητα ηλεκτρικής ενέργειας.} Η συχνότητα και οι τάσεις πρέπει να
      διατηρούνται εντός προδιαγεγραμμένων ορίων, τα οποία για την Ηπειρωτική Ευρώπη
      καθορίζονται από το κανονιστικό πλαίσιο.
\end{itemize}

Οι απαιτήσεις αυτές αφορούν φαινόμενα που εκτείνονται σε περισσότερες από δώδεκα τάξεις
μεγέθους χρόνου, από τα μικροδευτερόλεπτα έως τα έτη, και καθεμία εξυπηρετείται από
διαφορετικό μηχανισμό ελέγχου, όπως συνοψίζει ο Πίνακας~\ref{tab:chronikes-klimakes}.

\begin{table}[!htbp]
\centering
\begin{tabular}{lll}
\hline
Χρονική κλίμακα & Φαινόμενο ή διαδικασία & Μηχανισμός \\
\hline
μs -- ms   & βραχυκυκλώματα, υπερτάσεις   & προστασία, διακόπτες ισχύος \\
ms -- s    & ηλεκτρομηχανικές ταλαντώσεις & αδράνεια, PSS, FACTS \\
s -- 30 s  & απόκλιση συχνότητας          & εφεδρεία διατήρησης (ΕΔΣ) \\
30 s -- 15 min & αποκατάσταση συχνότητας  & εφεδρεία αποκατάστασης (ΕΑΣ) \\
min -- h   & αναπρογραμματισμός           & ενδοημερήσια αγορά \\
h -- ημέρα & ένταξη μονάδων, κατανομή     & αγορά επόμενης ημέρας \\
εβδομάδες -- έτη & επάρκεια ισχύος, συντηρήσεις & προγραμματισμός, επενδύσεις \\
\hline
\end{tabular}
\caption{Χρονικές κλίμακες λειτουργίας του ΣΜΗΕ και αντίστοιχοι μηχανισμοί ελέγχου. Το
κεφάλαιο εξετάζει κυρίως τις κλίμακες από τα δευτερόλεπτα και άνω.}
\label{tab:chronikes-klimakes}
\end{table}

\subsection{Καταστάσεις λειτουργίας και κριτήριο N-1}

Η κατάταξη της λειτουργίας σε διακριτές καταστάσεις ασφάλειας, με διαφορετική στρατηγική
ελέγχου σε καθεμία, διατυπώθηκε από τον Dy Liacco στα τέλη της δεκαετίας του 1960. Ο
Κανονισμός (ΕΕ) 2017/1485 για τη λειτουργία του συστήματος μεταφοράς (System Operation
Guideline, SO GL) ορίζει στο άρθρο 18 πέντε καταστάσεις του συστήματος:

\begin{enumerate}
\item \textbf{Κανονική κατάσταση} (normal state). Οι τάσεις, οι ροές ισχύος και η συχνότητα
      βρίσκονται εντός ορίων και το σύστημα παραμένει εντός ορίων μετά από οποιοδήποτε
      απρόβλεπτο συμβάν του καταλόγου απρόβλεπτων συμβάντων, λαμβανομένων υπόψη των
      διαθέσιμων διορθωτικών μέτρων. Ο έλεγχος είναι \emph{προληπτικός} (preventive).
\item \textbf{Κατάσταση συναγερμού} (alert state· συχνά και «κατάσταση επιφυλακής»). Τα
      μεγέθη βρίσκονται εντός ορίων, αλλά τουλάχιστον ένα απρόβλεπτο συμβάν του καταλόγου
      οδηγεί σε παραβίαση ορίων ακόμη και μετά την εφαρμογή διορθωτικών μέτρων, ή η
      εφεδρική δυναμικότητα έχει μειωθεί κατά περισσότερο από 20\% για διάστημα
      μεγαλύτερο των 30 λεπτών.
\item \textbf{Κατάσταση έκτακτης ανάγκης} (emergency state). Έχει παραβιαστεί τουλάχιστον
      ένα όριο επιχειρησιακής ασφάλειας ή εφαρμόζεται μέτρο του σχεδίου άμυνας του
      συστήματος. Ο έλεγχος είναι \emph{διορθωτικός} (corrective).
\item \textbf{Κατάσταση γενικής διακοπής} (blackout state· συχνά και «συσκότιση»). Έχει
      απολεσθεί περισσότερο από το 50\% της ζήτησης ή απουσιάζει πλήρως η τάση για
      τουλάχιστον τρία λεπτά στην περιοχή ελέγχου του ΔΣΜ.
\item \textbf{Κατάσταση αποκατάστασης} (restoration state). Ο ΔΣΜ έχει αρχίσει να
      εφαρμόζει τα μέτρα του σχεδίου αποκατάστασης για την επανατροφοδότηση του
      συστήματος.
\end{enumerate}

Η μετάβαση από την κανονική κατάσταση στην κατάσταση συναγερμού δεν συνοδεύεται από
παραβίαση ορίου ή από μεταβολή των μετρούμενων μεγεθών· προκύπτει μόνο από την ανάλυση
των απρόβλεπτων συμβάντων. Ο προσδιορισμός της κατάστασης του συστήματος προϋποθέτει
επομένως ακριβές μοντέλο της τρέχουσας λειτουργίας.

\subsubsection*{Το κριτήριο N-1}

Σύμφωνα με το \textbf{κριτήριο N-1}, η απώλεια οποιουδήποτε μεμονωμένου στοιχείου από τα $N$
στοιχεία του συστήματος — γραμμής, μετασχηματιστή, μονάδας ηλεκτροπαραγωγής ή ζυγού —
πρέπει να αντιμετωπίζεται χωρίς παραβίαση ορίων επιχειρησιακής ασφάλειας και χωρίς
αλυσιδωτές αποζεύξεις.

Η συμμόρφωση ελέγχεται με \emph{ανάλυση απρόβλεπτων συμβάντων} (contingency analysis· στη
βιβλιογραφία και «ανάλυση ενδεχομένων»): για κάθε συμβάν του καταλόγου εκτελείται
υπολογισμός ροής φορτίου (load flow) με το αντίστοιχο στοιχείο εκτός λειτουργίας και
ελέγχονται τα όρια. Για σύστημα μερικών εκατοντάδων κλάδων απαιτούνται εκατοντάδες
υπολογισμοί ροής φορτίου σε κάθε κύκλο ανάλυσης, ο οποίος πρέπει να ολοκληρώνεται εντός
λίγων λεπτών. Για τη μείωση του υπολογιστικού φόρτου εφαρμόζεται προκαταρκτική επιλογή των
κρισιμότερων συμβάντων (contingency screening), συνήθως με γραμμικοποιημένους συντελεστές
κατανομής (distribution factors).

Το κριτήριο N-1 είναι \emph{ντετερμινιστικό}: δεν λαμβάνει υπόψη την πιθανότητα εμφάνισης
κάθε συμβάντος και αντιμετωπίζει ισοδύναμα τα σπάνια και τα συχνά συμβάντα. Είναι απλό στην
εφαρμογή και στον έλεγχο, αλλά σε ορισμένες συνθήκες οδηγεί σε συντηρητική λειτουργία και σε
άλλες δεν καλύπτει πολλαπλά ή συσχετισμένα συμβάντα· για τον λόγο αυτόν αναπτύσσονται
πιθανοτικά κριτήρια αξιοπιστίας. Η ανάλυση απρόβλεπτων συμβάντων χρησιμοποιεί ως αρχική
συνθήκη την τρέχουσα κατάσταση του συστήματος, η οποία παρέχεται από την εκτίμηση
κατάστασης· σφάλμα της εκτίμησης μεταφέρεται απευθείας στην αξιολόγηση της ασφάλειας.

\subsection{Εποπτεία και έλεγχος: SCADA και EMS}

\subsubsection*{Αλυσίδα μέτρησης}

Η υποδομή τηλεμέτρησης και τηλεχειρισμού του ΣΜΗΕ βασίζεται στο σύστημα εποπτικού ελέγχου
και συλλογής δεδομένων (\textbf{SCADA}, Supervisory Control and Data Acquisition). Η αλυσίδα
από το φυσικό μέγεθος έως το κέντρο ελέγχου περιλαμβάνει τα ακόλουθα στοιχεία:

\begin{enumerate}
\item \textbf{Μετασχηματιστές μέτρησης.} Οι μετασχηματιστές τάσης (VT) και έντασης (CT)
      υποβιβάζουν τα μεγέθη του δικτύου σε τιμές κατάλληλες για τις διατάξεις μέτρησης
      και προστασίας, τυπικά 100 V και 1 ή 5 A στο δευτερεύον τύλιγμα. Εισάγουν το πρώτο
      σφάλμα της αλυσίδας, το οποίο καθορίζεται από την κλάση ακρίβειας (τυπικά 0,2 έως 1,
      δηλαδή σφάλμα λόγου 0,2 έως 1\%).
\item \textbf{Μετατροπείς μέτρησης (transducers) και ευφυείς ηλεκτρονικές συσκευές (IED).}
      Δειγματοληπτούν τις κυματομορφές και υπολογίζουν τα παράγωγα μεγέθη: ενεργό και
      άεργο ισχύ, ενεργό τιμή (rms) τάσης και έντασης.
\item \textbf{Απομακρυσμένες τερματικές μονάδες (RTU).} Συγκεντρώνουν τα δεδομένα του
      υποσταθμού, τα μεταδίδουν στο Κέντρο Ελέγχου Ενέργειας και εκτελούν τις εντολές
      τηλεχειρισμού.
\item \textbf{Τηλεπικοινωνιακό δίκτυο.} Μεταφέρει τα δεδομένα με τυποποιημένα πρωτόκολλα,
      κυρίως IEC 60870-5-101/104 και DNP3· εντός του υποσταθμού χρησιμοποιείται σήμερα
      κατά κανόνα το IEC 61850.
\end{enumerate}

Τα δεδομένα διακρίνονται σε δύο κατηγορίες:

\begin{itemize}
\item \textbf{Αναλογικές μετρήσεις} (τηλεμετρήσεις): ροές ενεργού και αέργου ισχύος στα
      άκρα των γραμμών και των μετασχηματιστών, εγχύσεις ισχύος στους ζυγούς, μέτρα τάσεων
      ζυγών, εντάσεις ρεύματος κλάδων.
\item \textbf{Ψηφιακές ενδείξεις} (τηλεσημάνσεις): θέσεις διακοπτών ισχύος και αποζευκτών,
      θέσεις των μεταγωγέων λήψεων υπό φορτίο των μετασχηματιστών, σήματα συναγερμού και
      λειτουργίας των ηλεκτρονόμων προστασίας. Οι ενδείξεις θέσης καθορίζουν την
      \emph{τοπολογία} του δικτύου, δηλαδή τα στοιχεία που βρίσκονται σε λειτουργία και τον
      τρόπο σύνδεσής τους.
\end{itemize}

Η περίοδος σάρωσης ενός συμβατικού συστήματος SCADA είναι τυπικά 2 έως 10 s.

\subsubsection*{Περιορισμοί των μετρήσεων SCADA}

Οι μετρήσεις SCADA παρουσιάζουν δύο περιορισμούς, οι οποίοι καθορίζουν τις απαιτήσεις της
εκτίμησης κατάστασης:

\begin{itemize}
\item \textbf{Απουσία χρονικού συγχρονισμού.} Οι μετρήσεις λαμβάνονται σε διαφορετικές
      χρονικές στιγμές εντός του κύκλου σάρωσης και δεν αντιστοιχούν σε κοινό χρονικό
      στιγμιότυπο. Το σύνολό τους αποτελεί \emph{οιονεί στατική} (quasi-static)
      αναπαράσταση του συστήματος, αποδεκτή σε μόνιμη κατάσταση λειτουργίας αλλά
      ανακριβή κατά τη διάρκεια μεταβατικών φαινομένων.
\item \textbf{Θόρυβος, εσφαλμένα δεδομένα και ελλιπής κάλυψη.} Οι μετρήσεις περιέχουν
      τυχαία σφάλματα και ενίοτε χονδροειδή σφάλματα (gross errors), ενώ ορισμένα μεγέθη
      δεν μετρώνται. Οι γωνίες των τάσεων των ζυγών, οι οποίες αποτελούν το ήμισυ των
      μεταβλητών κατάστασης, δεν μετρώνται από το συμβατικό SCADA.
\end{itemize}

Χονδροειδή σφάλματα προκαλούνται, μεταξύ άλλων, από εσφαλμένα καταχωρημένο λόγο
μετασχηματισμού μετασχηματιστή έντασης, από ανεστραμμένη πολικότητα μέτρησης ισχύος ή από
ένδειξη θέσης διακόπτη που δεν ανανεώνεται λόγω βλάβης. Τα δεδομένα αυτά δεν διακρίνονται
από τα ορθά με έλεγχο κάθε μέτρησης χωριστά· ανιχνεύονται μόνο με σύγκριση των μετρήσεων
μεταξύ τους, με αξιοποίηση του πλεονασμού των μετρήσεων (redundancy).

\subsubsection*{Λειτουργίες του EMS}

Το σύστημα διαχείρισης ενέργειας (\textbf{EMS}, Energy Management System) είναι το σύνολο των
εφαρμογών λογισμικού του κέντρου ελέγχου που επεξεργάζονται τα δεδομένα του SCADA. Οι
βασικές λειτουργίες εκτελούνται με την ακόλουθη σειρά, καθώς η έξοδος καθεμίας αποτελεί
είσοδο της επόμενης:

\begin{enumerate}
\item \textbf{Επεξεργαστής τοπολογίας} (topology processor). Προσδιορίζει, από τις θέσεις
      των διακοπτών και των αποζευκτών, την τρέχουσα τοπολογία και σχηματίζει το μοντέλο
      ζυγών-κλάδων (bus-branch model) του δικτύου.
\item \textbf{Εκτίμηση κατάστασης}. Προσδιορίζει τις τάσεις (μέτρα και γωνίες) όλων των
      ζυγών από τις διαθέσιμες μετρήσεις.
\item \textbf{Ανάλυση απρόβλεπτων συμβάντων}. Αξιολογεί την επιχειρησιακή ασφάλεια της
      εκτιμώμενης κατάστασης έναντι του καταλόγου απρόβλεπτων συμβάντων.
\item \textbf{Βέλτιστη ροή φορτίου} (Optimal Power Flow, OPF). Προσδιορίζει, εφόσον
      απαιτείται, τα αναγκαία διορθωτικά μέτρα (remedial actions) με το ελάχιστο κόστος.
\item \textbf{Αυτόματη Ρύθμιση Παραγωγής} (ΑΡΠ· Automatic Generation Control, AGC).
      Υλοποιεί συνεχώς τη ρύθμιση φορτίου-συχνότητας και την τήρηση των προγραμματισμένων
      ανταλλαγών ισχύος με τις γειτονικές περιοχές.
\end{enumerate}

Σφάλμα του επεξεργαστή τοπολογίας ή της εκτίμησης κατάστασης μεταφέρεται σε όλες τις
επόμενες λειτουργίες χωρίς να γίνεται αντιληπτό.

Το Σχήμα~\ref{fig:scada-ems} απεικονίζει τον κλειστό βρόχο εποπτείας και ελέγχου: οι
μετρήσεις μεταδίδονται από το δίκτυο προς το κέντρο ελέγχου και οι εντολές ελέγχου
επιστρέφουν προς το δίκτυο.

\begin{figure}[!htbp]
\centering
\includegraphics[width=\linewidth]{figures/scada-ems.svg}
\caption{Ο βρόχος εποπτείας και ελέγχου. Η επάνω σειρά είναι η αλυσίδα μέτρησης, από το
φυσικό δίκτυο έως το SCADA· η κάτω σειρά είναι η ακολουθία λειτουργιών του EMS, η οποία
καταλήγει σε εντολές προς το δίκτυο.}
\label{fig:scada-ems}
\end{figure}

\subsection{Συγχρονισμένες μετρήσεις φασιθετών και συστήματα WAMS}

\subsubsection*{Γωνία τάσης και ροή ενεργού ισχύος}

Η ενεργός ισχύς που μεταφέρεται μεταξύ δύο ζυγών μέσω γραμμής χωρίς απώλειες με επαγωγική
αντίδραση σειράς $X$ δίνεται από τη σχέση

\begin{equation}
P_{12} = \frac{V_1 V_2}{X}\,\sin(\delta_1 - \delta_2)
\label{eq:powerflow}
\end{equation}

Για τάσεις κοντά στην ονομαστική τιμή, η ροή ενεργού ισχύος καθορίζεται κυρίως από τη
διαφορά γωνιών $\delta_1 - \delta_2$, η οποία αποτελεί άμεσο μέτρο της φόρτισης της γραμμής
και του περιθωρίου στατικής ευστάθειας· η μέγιστη μεταφερόμενη ισχύς αντιστοιχεί σε
$\delta_1 - \delta_2 = 90^\circ$. Το συμβατικό SCADA δεν μετρά γωνίες τάσης.

\subsubsection*{Αρχή λειτουργίας της PMU}

Οι \textbf{μονάδες μέτρησης φασιθετών} (Phasor Measurement Units, PMU) υπολογίζουν τους
φασιθέτες — μέτρο και γωνία — της τάσης και της έντασης του ρεύματος, ανά φάση και για τη
θετική ακολουθία, με χρονοσήμανση ως προς την κοινή χρονική αναφορά UTC, η οποία λαμβάνεται
τυπικά από δέκτη GPS.

Η γωνία ενός φασιθέτη ορίζεται ως προς σήμα αναφοράς. Επειδή όλες οι PMU χρησιμοποιούν την
ίδια χρονική βάση, οι μετρούμενες γωνίες είναι άμεσα συγκρίσιμες σε ολόκληρο το σύστημα,
ανεξάρτητα από τη γεωγραφική απόσταση των σημείων μέτρησης· οι μετρήσεις αυτές ονομάζονται
\emph{συγχρονισμένοι φασιθέτες} (synchrophasors). Η απαιτούμενη ακρίβεια χρονισμού είναι της
τάξης του μικροδευτερολέπτου: σε σύστημα 50 Hz, σφάλμα χρονισμού 1 μs αντιστοιχεί σε σφάλμα
γωνίας $360^\circ \times 50 \times 10^{-6}$, δηλαδή 0,018°.

\begin{figure}[!htbp]
\centering
\includegraphics[width=\linewidth]{figures/pmu-wams.svg}
\caption{Αριστερά, το μετρούμενο μέγεθος: ο φασιθέτης τάσης με μέτρο και γωνία ως προς
κοινή χρονική αναφορά. Δεξιά, η ιεραρχία συλλογής: οι PMU συγχρονίζονται μέσω GPS, τα
πλαίσια δεδομένων ευθυγραμμίζονται χρονικά στον συγκεντρωτή δεδομένων (PDC) και
τροφοδοτούν το WAMS.}
\label{fig:pmu-wams}
\end{figure}

\subsubsection*{Προδιαγραφές}

Τα τεχνικά χαρακτηριστικά των PMU καθορίζονται από το πρότυπο IEC/IEEE 60255-118-1:2018, το
οποίο αντικατέστησε τη σειρά IEEE C37.118:

\begin{itemize}
\item \textbf{Ρυθμός αναφοράς} (reporting rate): τυπικά 10, 25 ή 50 πλαίσια δεδομένων ανά
      δευτερόλεπτο σε σύστημα 50 Hz· επιτρέπονται και υψηλότεροι ρυθμοί.
\item \textbf{Ακρίβεια}: το ολικό διανυσματικό σφάλμα (Total Vector Error, TVE) δεν
      επιτρέπεται να υπερβαίνει το 1\% σε μόνιμη κατάσταση. Το TVE ορίζεται ως το μέτρο της
      διαφοράς μετρούμενου και πραγματικού φασιθέτη, κανονικοποιημένο ως προς το μέτρο του
      πραγματικού, και συνδυάζει σε ένα μέγεθος τα σφάλματα μέτρου, γωνίας και χρονισμού.
      Ορίζονται επιπλέον το σφάλμα συχνότητας (FE) και το σφάλμα ρυθμού μεταβολής
      συχνότητας (RFE).
\item \textbf{Κλάσεις}: η κλάση P (protection) προορίζεται για εφαρμογές που απαιτούν ταχεία
      απόκριση και μικρή καθυστέρηση φίλτρου, ενώ η κλάση M (measurement) για εφαρμογές
      που απαιτούν υψηλή ακρίβεια παρουσία αρμονικών και παρεμβολών. Η διάκριση απορρέει
      από τον συμβιβασμό μεταξύ ταχύτητας απόκρισης και απόρριψης παρεμβολών του
      ψηφιακού φίλτρου.
\end{itemize}

\begin{table}[!htbp]
\centering
\begin{tabular}{lll}
\hline
 & SCADA & PMU / WAMS \\
\hline
Ανανέωση         & ανά 2--10 s        & 10--50 πλαίσια ανά s \\
Συγχρονισμός     & όχι                & ναι (UTC μέσω GPS) \\
Γωνία τάσης      & δεν μετράται       & μετράται άμεσα \\
Σχέση με την κατάσταση & μη γραμμική & γραμμική \\
Κάλυψη           & εκτεταμένη         & επιλεκτική, λόγω κόστους \\
Τυπική χρήση     & εκτίμηση κατάστασης, τηλεχειρισμοί & εποπτεία δυναμικών φαινομένων \\
\hline
\end{tabular}
\caption{Σύγκριση μετρήσεων SCADA και PMU. Οι δύο τεχνολογίες χρησιμοποιούνται
συμπληρωματικά.}
\label{tab:scada-pmu}
\end{table}

\subsubsection*{Συστήματα WAMS και εφαρμογές}

Οι μετρήσεις των PMU συγκεντρώνονται ιεραρχικά σε συγκεντρωτές δεδομένων φασιθετών (Phasor
Data Concentrators, PDC), οι οποίοι ευθυγραμμίζουν χρονικά τα πλαίσια δεδομένων και τα
διαθέτουν στο \textbf{σύστημα εποπτείας ευρείας περιοχής} (Wide Area Monitoring System,
WAMS). Οι κύριες εφαρμογές είναι:

\begin{itemize}
\item η παρακολούθηση ηλεκτρομηχανικών ταλαντώσεων μεταξύ περιοχών (inter-area
      oscillations), συχνότητας 0,1 έως 1 Hz, οι οποίες δεν αποτυπώνονται με περίοδο
      σάρωσης δευτερολέπτων·
\item η μέτρηση του ρυθμού μεταβολής της συχνότητας και η εκτίμηση της διαθέσιμης
      αδράνειας·
\item η ανίχνευση νησιδοποίησης, δηλαδή της μετάβασης τμήματος του δικτύου σε απομονωμένη
      λειτουργία, και η επιτήρηση των γωνιακών διαφορών μεταξύ περιοχών·
\item η επικύρωση δυναμικών μοντέλων μετά από διαταραχές, με σύγκριση της καταγεγραμμένης
      και της προσομοιωμένης απόκρισης·
\item η υποστήριξη της εκτίμησης κατάστασης και, σε συνθήκες πλήρους κάλυψης, η
      γραμμική εκτίμηση κατάστασης μόνο από μετρήσεις PMU.
\end{itemize}

\subsection{Ευέλικτα συστήματα μεταφοράς εναλλασσόμενου ρεύματος (FACTS)}

Οι διατάξεις \textbf{FACTS} (Flexible AC Transmission Systems) είναι διατάξεις ηλεκτρονικών
ισχύος που επιτρέπουν τον ταχύ και συνεχή έλεγχο των μεγεθών που καθορίζουν τη ροή ισχύος
σε δίκτυο εναλλασσόμενου ρεύματος.

Σύμφωνα με την \eqref{eq:powerflow}, η ροή ενεργού ισχύος εξαρτάται από τα μέτρα των τάσεων
$V_1, V_2$, την αντίδραση σειράς $X$ και τη διαφορά γωνιών $\delta_1 - \delta_2$. Σε
συμβατικό δίκτυο τα μεγέθη αυτά δεν ελέγχονται άμεσα, οπότε η κατανομή των ροών καθορίζεται
από τις σύνθετες αντιστάσεις των κλάδων και όχι από τις απαιτήσεις λειτουργίας. Οι διατάξεις
FACTS καθιστούν ελεγχόμενο ένα ή περισσότερα από τα μεγέθη αυτά και κατατάσσονται ανάλογα με
το μέγεθος που επηρεάζουν (Σχήμα~\ref{fig:facts}):

\begin{itemize}
\item \textbf{Παράλληλης σύνδεσης} (έλεγχος του $V$): ο στατός αντισταθμιστής αέργου ισχύος
      (Static Var Compensator, SVC) και ο στατός σύγχρονος αντισταθμιστής (Static
      Synchronous Compensator, STATCOM). Ρυθμίζουν την τάση του ζυγού μέσω έγχυσης ή
      απορρόφησης αέργου ισχύος.
\item \textbf{Σειριακής σύνδεσης} (έλεγχος του $X$): ο πυκνωτής σειράς ελεγχόμενος με
      θυρίστορ (Thyristor Controlled Series Capacitor, TCSC) και ο στατός σύγχρονος
      αντισταθμιστής σειράς (Static Synchronous Series Compensator, SSSC). Μεταβάλλουν την
      ισοδύναμη αντίδραση σειράς της γραμμής, ανακατανέμουν τις ροές ισχύος και αυξάνουν τη
      μεταφορική ικανότητα.
\item \textbf{Συνδυασμένης σύνδεσης}, σε σειρά και παράλληλα (έλεγχος και της γωνίας): ο
      ενοποιημένος ελεγκτής ροής ισχύος (Unified Power Flow Controller, UPFC) ελέγχει
      ταυτόχρονα και ανεξάρτητα την τάση και τις ροές ενεργού και αέργου ισχύος. Παρέχει τη
      μεγαλύτερη ευελιξία ελέγχου, με το υψηλότερο κόστος.
\end{itemize}

\begin{figure}[!htbp]
\centering
\includegraphics[width=0.88\linewidth]{figures/facts.svg}
\caption{Οι τρεις κατηγορίες διατάξεων FACTS σε γραμμή δύο ζυγών. Κάθε κατηγορία επεμβαίνει
σε διαφορετικό μέγεθος της σχέσης μεταφοράς ισχύος: η παράλληλη στην τάση, η σειριακή στην
αντίδραση, η συνδυασμένη και στη γωνία.}
\label{fig:facts}
\end{figure}

\subsubsection*{Τοπολογίες διατάξεων}

Ανάλογα με την αρχή λειτουργίας τους, οι διατάξεις FACTS διακρίνονται σε διατάξεις
\emph{μεταβλητής σύνθετης αντίστασης} (εμπέδησης), με θυρίστορ, και σε διατάξεις
\emph{ελεγχόμενης πηγής τάσης}, με μετατροπείς πηγής τάσης. Τα μονογραμμικά διαγράμματα
δίνονται στο Σχήμα~\ref{fig:facts-devices}.

\begin{itemize}
\item Ο \textbf{SVC} αποτελείται από πηνίο ελεγχόμενο με θυρίστορ (Thyristor Controlled
      Reactor, TCR) και πυκνωτές μεταγόμενους με θυρίστορ (Thyristor Switched Capacitor,
      TSC), σε παράλληλη σύνδεση. Η γωνία έναυσης του TCR ρυθμίζει συνεχώς την ισοδύναμη
      επαγωγική επιδεκτικότητα, ενώ οι TSC προσθέτουν χωρητική επιδεκτικότητα σε διακριτά
      βήματα. Η διάταξη συμπεριφέρεται ως \emph{μεταβλητή επιδεκτικότητα} $B$ (susceptance).
\item Ο \textbf{STATCOM} είναι μετατροπέας πηγής τάσης (Voltage Source Converter, VSC) με
      πυκνωτή στην πλευρά συνεχούς ρεύματος (ΣΡ). Παράγει εναλλασσόμενη τάση $V_\mu$
      ελεγχόμενου μέτρου, η οποία εφαρμόζεται στο δίκτυο μέσω της αντίδρασης σύζευξης: για
      $V_\mu > V$ ο STATCOM εγχέει άεργο ισχύ, ενώ για $V_\mu < V$ απορροφά.
\item Ο \textbf{TCSC} είναι πυκνωτής σειράς με παράλληλο TCR, οπότε η ισοδύναμη αντίδραση
      σειράς της γραμμής μεταβάλλεται συνεχώς. Ο \textbf{SSSC} εγχέει τάση σε σειρά μέσω
      μετασχηματιστή σύζευξης σειράς· η εγχεόμενη τάση είναι ανεξάρτητη του ρεύματος της
      γραμμής, ενώ στον TCSC η τάση στα άκρα του πυκνωτή είναι ανάλογη του ρεύματος.
\item Ο \textbf{UPFC} αποτελείται από έναν παράλληλο και έναν σειριακό VSC με \emph{κοινό
      ζυγό ΣΡ} (DC link). Ο κοινός ζυγός επιτρέπει την ανταλλαγή \emph{ενεργού} ισχύος
      μεταξύ των δύο μετατροπέων, οπότε ο σειριακός μετατροπέας εγχέει τάση οποιουδήποτε
      μέτρου και γωνίας εντός των ονομαστικών του ορίων και ελέγχει ανεξάρτητα τις ροές
      ενεργού και αέργου ισχύος.
\end{itemize}

\begin{figure}[!htbp]
\centering
\includegraphics[width=\linewidth]{figures/facts-devices.svg}
\caption{Μονογραμμικά διαγράμματα των βασικών διατάξεων FACTS. Οι δύο πρώτες συνδέονται
παράλληλα, οι δύο επόμενες σε σειρά, ενώ ο UPFC συνδυάζει αμφότερους τους τρόπους μέσω
κοινού ζυγού ΣΡ.}
\label{fig:facts-devices}
\end{figure}

\subsubsection*{Χαρακτηριστικές V--I των SVC και STATCOM}

Η διαφορά στην αρχή λειτουργίας αποτυπώνεται στις χαρακτηριστικές $V$--$I$ των δύο
διατάξεων παράλληλης σύνδεσης (Σχήμα~\ref{fig:facts-vi}).

Ο SVC συμπεριφέρεται ως μεταβλητή επιδεκτικότητα, οπότε $I = B\,V$ και $Q = B\,V^2$. Στο
όριο χωρητικής λειτουργίας το ρεύμα είναι ανάλογο της τάσης του ζυγού: στο 50\% της
ονομαστικής τάσης ο SVC αποδίδει το 50\% του ονομαστικού ρεύματος και το 25\% της ονομαστικής
αέργου ισχύος. Ο STATCOM, ως πηγή τάσης με έλεγχο ρεύματος, διατηρεί το ονομαστικό ρεύμα έως
πολύ χαμηλές τιμές τάσης, με περιορισμό μόνο από το μέγιστο επιτρεπόμενο ρεύμα των
ημιαγωγικών διακοπτών· η άεργος ισχύς του μειώνεται επομένως γραμμικά και όχι τετραγωνικά
με την τάση.

Κατά τη διάρκεια και αμέσως μετά από βραχυκύκλωμα, όταν η τάση είναι χαμηλή και η ανάγκη
στήριξης τάσης μέγιστη, ο STATCOM αποδίδει σημαντικά μεγαλύτερη άεργο ισχύ από SVC ίδιας
ονομαστικής ισχύος.

\begin{figure}[!htbp]
\centering
\includegraphics[width=\linewidth]{figures/facts-vi.svg}
\caption{Χαρακτηριστικά $V$--$I$ των δύο διατάξεων παράλληλης σύνδεσης. Τα όρια του SVC
είναι ευθείες σταθερής επιδεκτικότητας που διέρχονται από την αρχή, ενώ του STATCOM είναι
κατακόρυφες ευθείες σταθερού ρεύματος.}
\label{fig:facts-vi}
\end{figure}

Οι διατάξεις FACTS, μαζί με τα συστήματα συνεχούς ρεύματος υψηλής τάσης (HVDC), στα οποία η
ροή ισχύος καθορίζεται απευθείας από τον έλεγχο των μετατροπέων, αποτελούν τα κύρια μέσα
ρύθμισης τάσης και διαχείρισης συμφορήσεων στο σύστημα μεταφοράς. Τα αντίστοιχα μέσα στα
δίκτυα διανομής εξετάζονται στα Κεφάλαια 5 και 6.

\subsection{Εκτίμηση κατάστασης}

\subsubsection*{Διατύπωση του προβλήματος}

Η \textbf{εκτίμηση κατάστασης} (state estimation) προσδιορίζει την πλέον πιθανή κατάσταση
λειτουργίας του δικτύου από ένα πλεονάζον σύνολο μετρήσεων με θόρυβο, με βάση το μοντέλο
δικτύου που σχηματίζει ο επεξεργαστής τοπολογίας. Εισήχθη στα συστήματα ηλεκτρικής ενέργειας
από τους Schweppe και Wildes το 1970 και αποτελεί βασική λειτουργία κάθε EMS.

Ο υπολογισμός ροής φορτίου δεν είναι κατάλληλος για τον σκοπό αυτόν, για δύο λόγους.
Πρώτον, απαιτεί ακριβώς τόσα δεδομένα εισόδου όσοι είναι οι άγνωστοι, ενώ οι διαθέσιμες
μετρήσεις είναι περισσότερες και περιέχουν σφάλματα. Δεύτερον, δεν διαθέτει μηχανισμό
ανίχνευσης εσφαλμένων δεδομένων. Η εκτίμηση κατάστασης αξιοποιεί τον πλεονασμό των
μετρήσεων τόσο για τον περιορισμό της επίδρασης του θορύβου όσο και για την ανίχνευση και
την απόρριψη εσφαλμένων μετρήσεων.

Το Σχήμα~\ref{fig:se-flow} δίνει το διάγραμμα ροής του αλγορίθμου, ο οποίος περιλαμβάνει δύο
εμφωλευμένους βρόχους: τον εσωτερικό, για την επαναληπτική επίλυση του προβλήματος
ελαχιστοποίησης, και τον εξωτερικό, για την απόρριψη εσφαλμένων μετρήσεων.

\begin{figure}[!htbp]
\centering
\includegraphics[width=0.8\linewidth]{figures/se-flowchart.svg}
\caption{Ροή της εκτίμησης κατάστασης σε ένα EMS. Ο εσωτερικός βρόχος (αριστερά) επαναλαμβάνει
τη μέθοδο Gauss--Newton έως τη σύγκλιση· ο εξωτερικός (δεξιά) απορρίπτει τη μέτρηση με το
μέγιστο κανονικοποιημένο υπόλοιπο και επανεκτιμά.}
\label{fig:se-flow}
\end{figure}

\subsubsection*{Μοντέλο μετρήσεων}

Ως \emph{διάνυσμα κατάστασης} ορίζεται το διάνυσμα
$\mathbf{x} = [\boldsymbol{\theta}^{T}, \mathbf{V}^{T}]^{T}$ των γωνιών και των μέτρων των
τάσεων όλων των ζυγών. Για δίκτυο $N$ ζυγών, η γωνία του ζυγού αναφοράς (reference bus)
τίθεται ίση με μηδέν, οπότε το πλήθος των μεταβλητών κατάστασης είναι

\begin{equation}
n = 2N - 1
\label{eq:nstates}
\end{equation}

Από τα $\mathbf{V}$ και $\boldsymbol{\theta}$ υπολογίζεται κάθε άλλο ηλεκτρικό μέγεθος του
δικτύου (ροές, εγχύσεις, απώλειες) μέσω των εξισώσεων ροής φορτίου· το διάνυσμα κατάστασης
αποτελεί επομένως το ελάχιστο σύνολο μεταβλητών που περιγράφει πλήρως τη μόνιμη κατάσταση
λειτουργίας.

Το μοντέλο μετρήσεων είναι

\begin{equation}
\mathbf{z} = \mathbf{h}(\mathbf{x}) + \mathbf{e}
\label{eq:measmodel}
\end{equation}

όπου $\mathbf{z} \in \mathbb{R}^{m}$ το διάνυσμα των μετρήσεων, $\mathbf{h}(\cdot)$ το
διάνυσμα των μη γραμμικών συναρτήσεων που συνδέουν την κατάσταση με τα μετρούμενα μεγέθη και
$\mathbf{e}$ το διάνυσμα των σφαλμάτων μέτρησης. Τα σφάλματα θεωρούνται ανεξάρτητα, κανονικά
κατανεμημένα, μηδενικής μέσης τιμής, με διαγώνιο πίνακα συνδιακύμανσης
$\mathbf{R} = \mathrm{diag}(\sigma_1^2, \ldots, \sigma_m^2)$.

Ο \textbf{λόγος πλεονασμού} ορίζεται ως

\begin{equation}
\eta = \frac{m}{n}
\label{eq:redundancy}
\end{equation}

και σε πραγματικά συστήματα κυμαίνεται τυπικά μεταξύ 1,7 και 3. Για παράδειγμα, σε δίκτυο
$N = 100$ ζυγών είναι $n = 199$, οπότε για $\eta = 2$ απαιτούνται περίπου $m = 400$
μετρήσεις. Για $\eta = 1$ το πρόβλημα ανάγεται σε υπολογισμό ροής φορτίου και η ανίχνευση
εσφαλμένων δεδομένων δεν είναι δυνατή.

\subsubsection*{Σταθμισμένα ελάχιστα τετράγωνα}

Για $m > n$ το σύστημα \eqref{eq:measmodel} είναι υπερκαθορισμένο και δεν υπάρχει γενικά
$\mathbf{x}$ που ικανοποιεί όλες τις μετρήσεις. Η εκτίμηση προκύπτει από την ελαχιστοποίηση
της αντικειμενικής συνάρτησης των \textbf{σταθμισμένων ελαχίστων τετραγώνων} (Weighted Least
Squares, WLS):

\begin{equation}
J(\mathbf{x}) = \sum_{i=1}^{m} \frac{\left(z_i - h_i(\mathbf{x})\right)^2}{\sigma_i^2}
             = \left[\mathbf{z}-\mathbf{h}(\mathbf{x})\right]^{T} \mathbf{R}^{-1}
               \left[\mathbf{z}-\mathbf{h}(\mathbf{x})\right]
\label{eq:wls}
\end{equation}

Ο συντελεστής στάθμισης $1/\sigma_i^2$ αποδίδει μεγαλύτερο βάρος στις ακριβέστερες μετρήσεις.
Κάθε όρος του αθροίσματος είναι το τετράγωνο του υπολοίπου εκφρασμένου σε μονάδες τυπικής
απόκλισης, ώστε μετρήσεις διαφορετικών μεγεθών και κλιμάκων να είναι συγκρίσιμες. Για
κανονικά κατανεμημένα σφάλματα, η ελαχιστοποίηση της \eqref{eq:wls} ταυτίζεται με την
εκτίμηση μέγιστης πιθανοφάνειας.

Επειδή η $\mathbf{h}$ είναι μη γραμμική, η ελαχιστοποίηση γίνεται επαναληπτικά με τη μέθοδο
Gauss--Newton. Με $\mathbf{H}(\mathbf{x}) = \partial \mathbf{h}/\partial \mathbf{x}$ τον
ιακωβιανό πίνακα, σε κάθε επανάληψη $k$ επιλύονται οι \emph{κανονικές εξισώσεις}

\begin{align}
\mathbf{G}(\mathbf{x}^{k}) \, \Delta \mathbf{x}^{k}
  &= \mathbf{H}^{T}(\mathbf{x}^{k}) \, \mathbf{R}^{-1}
     \left[\mathbf{z} - \mathbf{h}(\mathbf{x}^{k})\right] \\
\mathbf{G}(\mathbf{x}^{k})
  &= \mathbf{H}^{T}(\mathbf{x}^{k}) \, \mathbf{R}^{-1} \mathbf{H}(\mathbf{x}^{k})
\label{eq:gain}
\end{align}

και η εκτίμηση ενημερώνεται ως $\mathbf{x}^{k+1} = \mathbf{x}^{k} + \Delta \mathbf{x}^{k}$,
έως ότου το $\max|\Delta \mathbf{x}^{k}|$ γίνει μικρότερο από προκαθορισμένη ανοχή
σύγκλισης. Με αρχική εκτίμηση τη λύση της προηγούμενης εκτέλεσης, η σύγκλιση επιτυγχάνεται
τυπικά σε τρεις έως πέντε επαναλήψεις.

Ο πίνακας $\mathbf{G}$ ονομάζεται \emph{πίνακας κέρδους} (gain matrix). Είναι συμμετρικός,
θετικά ορισμένος όταν το σύστημα είναι παρατηρήσιμο, και \emph{αραιός}, επειδή κάθε μέτρηση
συνδέεται μόνο με τις μεταβλητές κατάστασης γειτονικών ζυγών. Οι κανονικές εξισώσεις
επιλύονται με αραιή παραγοντοποίηση Cholesky, της οποίας το υπολογιστικό κόστος αυξάνει
σχεδόν γραμμικά με το μέγεθος του δικτύου, έναντι $O(n^3)$ για πυκνή παραγοντοποίηση. Η
αραιότητα καθιστά επομένως δυνατή την εκτέλεση σε πραγματικό χρόνο για δίκτυα χιλιάδων
ζυγών.

\subsubsection*{Παρατηρησιμότητα}

Η επίλυση των κανονικών εξισώσεων \eqref{eq:gain} προϋποθέτει ότι ο $\mathbf{G}$ είναι
αντιστρέψιμος,
δηλαδή ότι ο $\mathbf{H}$ έχει πλήρη τάξη στηλών $n$. Η συνθήκη εξαρτάται από τη θέση και το
είδος των μετρήσεων και όχι μόνο από το πλήθος τους: μετρήσεις συγκεντρωμένες σε τμήμα του
δικτύου δεν επιτρέπουν τον προσδιορισμό της κατάστασης στο υπόλοιπο δίκτυο.

Το σύστημα χαρακτηρίζεται \textbf{παρατηρήσιμο} όταν οι διαθέσιμες μετρήσεις επιτρέπουν τον
μονοσήμαντο προσδιορισμό ολόκληρου του διανύσματος κατάστασης. Η παρατηρησιμότητα ελέγχεται
με δύο μεθόδους:

\begin{itemize}
\item \textbf{Τοπολογική}: αναζητείται επικαλύπτον δένδρο (spanning tree) του γράφου του
      δικτύου, κάθε κλάδος του οποίου αντιστοιχίζεται σε μέτρηση. Βασίζεται στη θεωρία
      γράφων και δεν επηρεάζεται από σφάλματα στρογγυλοποίησης.
\item \textbf{Αριθμητική}: ελέγχεται η τάξη (rank) του $\mathbf{H}$ μέσω της τριγωνικής
      παραγοντοποίησης του $\mathbf{G}$, με εντοπισμό μηδενικών οδηγών (pivots). Είναι
      γενικότερη και αντιμετωπίζει και τις οριακές περιπτώσεις.
\end{itemize}

Σε μη παρατηρήσιμο σύστημα προσδιορίζονται οι \emph{παρατηρήσιμες νησίδες} (observable
islands) και η κατάσταση εκτιμάται εντός αυτών. Η παρατηρησιμότητα αποκαθίσταται με
\emph{ψευδομετρήσεις} (pseudo-measurements), δηλαδή εκτιμήσεις φορτίου από ιστορικά δεδομένα
ή προβλέψεις, στις οποίες αποδίδεται μεγάλη τυπική απόκλιση ώστε η επίδρασή τους στην
εκτίμηση να είναι μικρή.

\subsubsection*{Ανίχνευση και αναγνώριση εσφαλμένων μετρήσεων}

Η ανίχνευση εσφαλμένων δεδομένων (bad data) βασίζεται στα \emph{υπόλοιπα} (residuals)
$\mathbf{r} = \mathbf{z}-\mathbf{h}(\hat{\mathbf{x}})$ της εκτίμησης. Η κλασική διαδικασία
περιλαμβάνει δύο στάδια.

\textbf{Στάδιο 1 --- Ανίχνευση.} Όταν δεν υπάρχουν εσφαλμένα δεδομένα, η τιμή
$J(\hat{\mathbf{x}})$ ακολουθεί προσεγγιστικά κατανομή $\chi^2$ με $m-n$ βαθμούς ελευθερίας
και μέση τιμή $m-n$· για το προηγούμενο παράδειγμα ($m = 400$, $n = 199$) η αναμενόμενη τιμή
είναι 201. Αν $J(\hat{\mathbf{x}}) > \chi^2_{m-n,\,p}$, όπου $p$ η πιθανότητα του ελέγχου
(συνήθως 95\%), θεωρείται ότι υπάρχουν εσφαλμένα δεδομένα. Ο έλεγχος $\chi^2$ ανιχνεύει την
ύπαρξη εσφαλμένων δεδομένων, αλλά δεν τα εντοπίζει.

\textbf{Στάδιο 2 --- Αναγνώριση.} Τα υπόλοιπα έχουν διαφορετικές διασπορές και δεν
συγκρίνονται απευθείας. Κανονικοποιούνται ως

\begin{equation}
r_i^{N} = \frac{|r_i|}{\sqrt{\Omega_{ii}}}, \qquad
\boldsymbol{\Omega} = \mathbf{R} - \mathbf{H}\mathbf{G}^{-1}\mathbf{H}^{T}
\label{eq:normres}
\end{equation}

όπου $\boldsymbol{\Omega}$ ο πίνακας συνδιακύμανσης των υπολοίπων. Κατά το κριτήριο του
μεγίστου κανονικοποιημένου υπολοίπου (Largest Normalized Residual, LNR), αν το μέγιστο
$r_i^{N}$ υπερβαίνει ένα κατώφλι, συνήθως 3, η αντίστοιχη μέτρηση απορρίπτεται και η
εκτίμηση επαναλαμβάνεται. Η διαδικασία συνεχίζεται μέχρι να μην υπάρχει κανονικοποιημένο
υπόλοιπο μεγαλύτερο από το κατώφλι. Το κριτήριο LNR είναι αξιόπιστο για μεμονωμένα
εσφαλμένα δεδομένα, αλλά μπορεί να αστοχήσει όταν υπάρχουν πολλαπλά εσφαλμένα δεδομένα σε
μετρήσεις με ισχυρά συσχετισμένα υπόλοιπα.

\subsubsection*{Κρίσιμες μετρήσεις και επιθέσεις έγχυσης ψευδών δεδομένων}

Η ανίχνευση εσφαλμένων δεδομένων μέσω των υπολοίπων έχει δύο εγγενείς περιορισμούς.

\textbf{Κρίσιμες μετρήσεις.} Μια μέτρηση χαρακτηρίζεται \emph{κρίσιμη} (critical measurement)
όταν η αφαίρεσή της καθιστά το σύστημα μη παρατηρήσιμο. Για κρίσιμη μέτρηση ισχύει
$\Omega_{ii} = 0$ και το υπόλοιπό της είναι πάντοτε μηδενικό, επειδή καμία άλλη μέτρηση δεν
παρέχει πληροφορία για το αντίστοιχο μέγεθος. Σφάλμα σε κρίσιμη μέτρηση είναι επομένως μη
ανιχνεύσιμο, ανεξάρτητα από το μέγεθός του. Η αποφυγή κρίσιμων μετρήσεων αποτελεί κριτήριο
σχεδιασμού του συστήματος μετρήσεων.

\textbf{Επιθέσεις έγχυσης ψευδών δεδομένων.} Οι Liu, Ning και Reiter έδειξαν το 2009 ότι
επιτιθέμενος που γνωρίζει την τοπολογία και τις παραμέτρους του δικτύου, δηλαδή τον πίνακα
$\mathbf{H}$ του γραμμικού (DC) μοντέλου, μπορεί να αλλοιώσει τις μετρήσεις χωρίς να
ανιχνευθεί. Αν οι μετρήσεις αλλοιωθούν κατά $\mathbf{a} = \mathbf{H}\mathbf{c}$ για τυχόν
διάνυσμα $\mathbf{c}$, το διάνυσμα $\mathbf{z}_{\text{bad}} = \mathbf{z} + \mathbf{a}$ οδηγεί
στην εκτίμηση $\hat{\mathbf{x}} + \mathbf{c}$, ενώ τα υπόλοιπα παραμένουν αμετάβλητα:

\begin{equation}
\mathbf{r}_{\text{bad}} = \mathbf{z} + \mathbf{H}\mathbf{c}
        - \mathbf{H}(\hat{\mathbf{x}} + \mathbf{c}) = \mathbf{r}
\label{eq:fdia}
\end{equation}

Επειδή το διάνυσμα επίθεσης είναι συνεπές με το μοντέλο μετρήσεων, η επίθεση δεν
ανιχνεύεται από κανέναν έλεγχο που βασίζεται στα υπόλοιπα, όπως ο έλεγχος $\chi^2$ και το
κριτήριο LNR. Οι επιθέσεις έγχυσης ψευδών δεδομένων (False Data Injection Attacks, FDIA) και
τα μέτρα αντιμετώπισής τους εξετάζονται στο Κεφάλαιο 8.

\subsubsection*{Ενσωμάτωση μετρήσεων PMU}

Οι PMU μετρούν απευθείας φασιθέτες τάσης και ρεύματος. Σε καρτεσιανές (ορθογώνιες)
συντεταγμένες, οι φασιθέτες τάσης των ζυγών και ρεύματος των κλάδων είναι γραμμικές
συναρτήσεις των μιγαδικών τάσεων των ζυγών. Αν το σύστημα είναι πλήρως παρατηρήσιμο μόνο από
PMU, το μοντέλο \eqref{eq:measmodel} γίνεται γραμμικό,
$\mathbf{z} = \mathbf{H}\mathbf{x} + \mathbf{e}$, και η εκτίμηση υπολογίζεται χωρίς
επαναλήψεις:

\begin{equation}
\hat{\mathbf{x}} = \left(\mathbf{H}^{T}\mathbf{R}^{-1}\mathbf{H}\right)^{-1}
                    \mathbf{H}^{T}\mathbf{R}^{-1}\mathbf{z}
\label{eq:linse}
\end{equation}

Η γραμμική εκτίμηση δεν παρουσιάζει προβλήματα σύγκλισης ούτε εξάρτηση από την αρχική
εκτίμηση. Ο $\mathbf{H}$ είναι σταθερός για δεδομένη τοπολογία, οπότε ο πίνακας κέρδους
παραγοντοποιείται μία φορά, και η εκτίμηση μπορεί να εκτελείται με τον ρυθμό αναφοράς των
PMU, δηλαδή δεκάδες φορές ανά δευτερόλεπτο, αντί για μία φορά ανά ένα έως λίγα λεπτά.

Επειδή η πλήρης κάλυψη με PMU δεν είναι συνήθως οικονομικά εφικτή, χρησιμοποιούνται
\textbf{υβριδικοί εκτιμητές} που συνδυάζουν μετρήσεις SCADA και PMU. Η κύρια δυσκολία τους
είναι ο συνδυασμός μετρήσεων με διαφορετική ακρίβεια και διαφορετικό ρυθμό ανανέωσης,
ενδεικτικά 2 έως 10 s για το SCADA έναντι 20 ms για τις PMU.

\subsection{Στάδια λειτουργίας: από τον προγραμματισμό στον πραγματικό χρόνο}

Ο προγραμματισμός της λειτουργίας γίνεται σε διαδοχικούς χρονικούς ορίζοντες. Καθώς
πλησιάζει ο χρόνος φυσικής παράδοσης, η αβεβαιότητα των προβλέψεων φορτίου και ανανεώσιμης
παραγωγής μειώνεται, ενώ οι διαθέσιμες ενέργειες διόρθωσης περιορίζονται και το κόστος τους
αυξάνεται. Κάθε στάδιο διορθώνει τις αποκλίσεις του προηγούμενου. Στο ευρωπαϊκό
Μοντέλο-Στόχο (Target Model) οι ορίζοντες αυτοί αντιστοιχούν σε διακριτές αγορές:

\begin{itemize}
\item \textbf{Μακροπρόθεσμος και προθεσμιακός ορίζοντας} (έτη έως εβδομάδες). Μελέτες
      επάρκειας ισχύος, συντονισμός των προγραμμάτων συντήρησης γραμμών και μονάδων,
      διαστασιολόγηση των εφεδρειών, συναλλαγές στην Προθεσμιακή Αγορά για αντιστάθμιση
      του κινδύνου τιμής.
\item \textbf{Αγορά Επόμενης Ημέρας} (day-ahead). Αποτελεί το κύριο στάδιο
      προγραμματισμού. Οι συμμετέχοντες υποβάλλουν εντολές αγοράς και πώλησης για κάθε
      αγοραία χρονική μονάδα της επόμενης ημέρας. Η αγορά εκκαθαρίζεται μέσω της ενιαίας
      σύζευξης επόμενης ημέρας σε ευρωπαϊκό επίπεδο (Single Day-Ahead Coupling, SDAC), με
      μεγιστοποίηση του συνολικού κοινωνικού πλεονάσματος υπό τους περιορισμούς της
      διαθέσιμης διαζωνικής δυναμικότητας· οι ροές μεταξύ ζωνών προσφοράς κατευθύνονται
      από ζώνες χαμηλότερης προς ζώνες υψηλότερης τιμής. Ο ΔΣΜ ελέγχει στη συνέχεια την
      τεχνική εφικτότητα του προγράμματος.
\item \textbf{Ενδοημερήσια Αγορά} (intra-day). Επιτρέπει στους συμμετέχοντες να
      αναπροσαρμόζουν τις θέσεις τους καθώς βελτιώνονται οι προβλέψεις. Η σημασία της
      αυξάνεται με τη διείσδυση αιολικής και φωτοβολταϊκής παραγωγής, επειδή το σφάλμα
      πρόβλεψης της παραγωγής αυτής μειώνεται σημαντικά με τη μείωση του ορίζοντα
      πρόβλεψης.
\item \textbf{Πραγματικός χρόνος και εξισορρόπηση} (balancing). Ο ΔΣΜ ενεργοποιεί προσφορές
      ενέργειας εξισορρόπησης για την κάλυψη των εναπομενουσών ανισορροπιών, την
      αντιμετώπιση διαταραχών και τη διαχείριση συμφορήσεων.
\end{itemize}

\subsubsection*{Αγοραία χρονική μονάδα}

Η \emph{αγοραία χρονική μονάδα} (Market Time Unit, MTU) είναι η χρονική ανάλυση εκκαθάρισης
της αγοράς. Από την ημέρα παράδοσης 1η Οκτωβρίου 2025 (δημοπρασία της 30ής Σεπτεμβρίου
2025), η ενιαία σύζευξη επόμενης ημέρας λειτουργεί με αγοραία χρονική μονάδα \textbf{15
λεπτών} αντί της ωριαίας, σε εναρμόνιση με την περίοδο εκκαθάρισης ανισορροπιών. Με
μικρότερη αγοραία χρονική μονάδα, οι ταχείες μεταβολές φορτίου και παραγωγής, όπως οι
κλίσεις της φωτοβολταϊκής παραγωγής, αποτυπώνονται στο πρόγραμμα αντί να καλύπτονται με
ενέργεια εξισορρόπησης· μειώνονται έτσι οι αποκλίσεις στις αλλαγές ωριαίων περιόδων και οι
ανάγκες σε εφεδρείες.

\subsubsection*{Ελληνική αγορά ηλεκτρικής ενέργειας}

Στην Ελλάδα το Μοντέλο-Στόχος εφαρμόζεται από την 1η Νοεμβρίου 2020. Τη διαχείριση της
Αγοράς Επόμενης Ημέρας και της Ενδοημερήσιας Αγοράς έχει το Ελληνικό Χρηματιστήριο Ενέργειας
(ΕΧΕ) και τη διαχείριση της \textbf{Αγοράς Εξισορρόπησης} ο ΑΔΜΗΕ, ο οποίος εκτελεί τη
Διαδικασία Ολοκληρωμένου Προγραμματισμού και ενεργοποιεί σε πραγματικό χρόνο ενέργεια
εξισορρόπησης και εφεδρείες συχνότητας.

\subsection{Απόκριση συχνότητας και εφεδρείες}

\subsubsection*{Αδράνεια και αρχικός ρυθμός μεταβολής συχνότητας}

Σε μόνιμη κατάσταση η συχνότητα είναι σταθερή και κοινή σε ολόκληρη τη συγχρονισμένη
περιοχή· κάθε ανισορροπία μεταξύ παραγωγής και ζήτησης προκαλεί μεταβολή της.

Αμέσως μετά από ανισορροπία ισχύος $\Delta P$, για παράδειγμα την απώλεια μεγάλης μονάδας
ηλεκτροπαραγωγής, οι ρυθμιστές στροφών δεν έχουν ακόμη αποκριθεί και το έλλειμμα καλύπτεται
από την κινητική ενέργεια των στρεφόμενων μαζών των σύγχρονων μηχανών. Ο αρχικός ρυθμός
μεταβολής της συχνότητας (Rate of Change of Frequency, RoCoF) προκύπτει από την εξίσωση
ταλάντωσης (swing equation) και δίνεται κατ' απόλυτη τιμή από τη σχέση

\begin{equation}
\frac{\mathrm{d}f}{\mathrm{d}t} = \frac{f_0 \, \Delta P}{2 H S_n}
\label{eq:rocof}
\end{equation}

όπου $f_0$ η ονομαστική συχνότητα, $H$ η σταθερά αδράνειας του συστήματος σε δευτερόλεπτα
και $S_n$ η συνολική ονομαστική φαινόμενη ισχύς των \emph{συγχρονισμένων} μονάδων.

\emph{Παράδειγμα.} Για συγχρονισμένη περιοχή με $S_n = 300$ GVA, $H = 5$ s και απώλεια
$\Delta P = 3\,000$ MW, ο αρχικός ρυθμός μεταβολής είναι
$50 \times 3000 / (2 \times 5 \times 300\,000)$ = 0,05 Hz/s. Για την ίδια απώλεια σε
συνθήκες χαμηλού φορτίου, με το ήμισυ της συγχρονισμένης ισχύος, ο ρυθμός διπλασιάζεται.

Σύμφωνα με την \eqref{eq:rocof}, η αντικατάσταση σύγχρονων γεννητριών από μονάδες που
συνδέονται μέσω ηλεκτρονικών μετατροπέων ισχύος, όπως τα αιολικά και φωτοβολταϊκά πάρκα,
μειώνει την αποθηκευμένη κινητική ενέργεια $H S_n$ και αυξάνει τον ρυθμό μεταβολής της
συχνότητας. Μειώνεται έτσι ο διαθέσιμος χρόνος πριν η συχνότητα φθάσει στο κατώφλι της
αυτόματης αποσύνδεσης ζήτησης λόγω χαμηλής συχνότητας, και απαιτούνται ταχύτερες υπηρεσίες,
όπως η συνθετική αδράνεια και η ταχεία απόκριση συχνότητας.

Το Σχήμα~\ref{fig:freq} απεικονίζει την απόκριση συχνότητας για το συμβάν αναφοράς και τις
τρεις διαδοχικές φάσεις της.

\begin{figure}[!htbp]
\centering
\includegraphics[width=\linewidth]{figures/frequency-response.svg}
\caption{Απόκριση συχνότητας σε απώλεια παραγωγής 3\,000 MW. Η αρχική κλίση καθορίζεται
αποκλειστικά από την αδράνεια· η εφεδρεία διατήρησης συχνότητας (ΕΔΣ) ανακόπτει την πτώση
και σταθεροποιεί τη συχνότητα σε νέα τιμή μόνιμης κατάστασης· η εφεδρεία αποκατάστασης
συχνότητας (ΕΑΣ) εξαλείφει το μόνιμο σφάλμα και επαναφέρει τη συχνότητα στα 50 Hz.}
\label{fig:freq}
\end{figure}

\subsubsection*{Εφεδρεία διατήρησης συχνότητας και στατισμός}

Η \textbf{εφεδρεία διατήρησης συχνότητας} (ΕΔΣ· Frequency Containment Reserve, FCR)
υλοποιείται από τους ρυθμιστές στροφών των μονάδων, οι οποίοι μεταβάλλουν την παραγόμενη
ισχύ αναλογικά προς την απόκλιση συχνότητας, σύμφωνα με τον \emph{στατισμό} $R$:

\begin{equation}
\frac{\Delta P}{P_n} = -\frac{1}{R}\,\frac{\Delta f}{f_0}
\label{eq:droop}
\end{equation}

Τυπικές τιμές στατισμού είναι $R = 4$ έως 5\%. Για παράδειγμα, μονάδα 500 MW με $R = 5\%$,
για απόκλιση $\Delta f = -200$ mHz, δηλαδή −0,4\%, αυξάνει την παραγωγή της κατά
(0,004~/~0,05)~×~500~=~40~MW.

Το αρνητικό πρόσημο αντιστοιχεί σε αρνητική ανάδραση: μείωση της συχνότητας προκαλεί αύξηση
της παραγωγής. Επειδή ο έλεγχος είναι \emph{αναλογικός}, η πρόσθετη παραγωγή διατηρείται μόνο
όσο διατηρείται η απόκλιση συχνότητας. Για ανισορροπία $\Delta P$ η απόκλιση μόνιμης
κατάστασης είναι $\Delta f_{\mathrm{ss}} = -\Delta P / \lambda$, όπου $\lambda$ (MW/Hz) η
χαρακτηριστική ισχύος-συχνότητας της περιοχής, δηλαδή το άθροισμα της συνεισφοράς των
ρυθμιστών και της αυτορρύθμισης του φορτίου. Η ΕΔΣ επομένως σταθεροποιεί τη συχνότητα σε
νέα τιμή μόνιμης κατάστασης, με μόνιμο σφάλμα ως προς την ονομαστική τιμή, το οποίο
εξαλείφεται από την εφεδρεία αποκατάστασης συχνότητας.

\subsubsection*{Εφεδρεία αποκατάστασης συχνότητας σε απομονωμένο σύστημα}

Η \textbf{εφεδρεία αποκατάστασης συχνότητας} (ΕΑΣ· Frequency Restoration Reserve, FRR)
εξετάζεται αρχικά για \emph{απομονωμένο} σύστημα, χωρίς διασυνδέσεις με γειτονικά
συστήματα. Στην περίπτωση αυτή το μόνο μέγεθος ελέγχου είναι η απόκλιση συχνότητας και ο
κεντρικός ελεγκτής εφαρμόζει \emph{ολοκληρωτική} δράση:

\begin{equation}
\Delta P_{\mathrm{FRR}}(t) = -K_I \int_0^{t} \Delta f(\tau)\,\mathrm{d}\tau
\label{eq:frr-isolated}
\end{equation}

Όσο $\Delta f \neq 0$, το ολοκλήρωμα μεταβάλλεται και η εντολή παραγωγής αυξάνεται· σε μόνιμη
κατάσταση ισχύει κατ' ανάγκη $\Delta f = 0$. Ο ολοκληρωτικός έλεγχος εξαλείφει επομένως το
μόνιμο σφάλμα της αναλογικής ΕΔΣ. Με την επαναφορά της συχνότητας στην ονομαστική τιμή, η
συνεισφορά της ΕΔΣ μηδενίζεται σύμφωνα με την \eqref{eq:droop}, δηλαδή η ΕΑΣ
\emph{αποδεσμεύει} την ΕΔΣ, η οποία καθίσταται εκ νέου διαθέσιμη για επόμενη διαταραχή.

Η απόκριση της ΕΑΣ σχεδιάζεται σημαντικά βραδύτερη από εκείνη της ΕΔΣ (λεπτά έναντι
δευτερολέπτων). Ο διαχωρισμός των χρονικών κλιμάκων αποτρέπει την αλληλεπίδραση των δύο
βρόχων ελέγχου και τις ταλαντώσεις που θα προέκυπταν από αυτήν.

\subsubsection*{Εφεδρεία αποκατάστασης συχνότητας σε σύστημα πολλών περιοχών}

Σε διασυνδεδεμένο σύστημα η συχνότητα είναι κοινή σε ολόκληρη τη συγχρονισμένη περιοχή,
οπότε απώλεια παραγωγής σε οποιαδήποτε περιοχή προκαλεί απόκλιση συχνότητας σε όλες. Αν ο
ελεγκτής κάθε περιοχής χρησιμοποιούσε μόνο την απόκλιση συχνότητας, όλες οι περιοχές θα
ενεργοποιούσαν ΕΑΣ για την ίδια διαταραχή. Αυτό θα προκαλούσε υπερδιόρθωση, αποκλίσεις των
ανταλλαγών ισχύος από τα προγράμματα και μετακύλιση του κόστους εξισορρόπησης στις περιοχές
που δεν προκάλεσαν τη διαταραχή.

Για τον διαχωρισμό της περιοχής στην οποία εμφανίστηκε η ανισορροπία, ο ελεγκτής κάθε
περιοχής ελέγχου φορτίου-συχνότητας (περιοχή ΕΦΣ· LFC area) μηδενίζει το \emph{σφάλμα ελέγχου
περιοχής} (ΣΕΠ· Area Control Error, ACE):

\begin{equation}
\mathrm{ACE} = \Delta P_{\mathrm{tie}} + K \, \Delta f
\label{eq:ace}
\end{equation}

όπου $\Delta P_{\mathrm{tie}}$ το σφάλμα ελέγχου ισχύος, δηλαδή η απόκλιση της μετρούμενης
ανταλλαγής ισχύος στις γραμμές διασύνδεσης από το πρόγραμμα, και $K$ ο συντελεστής K της
περιοχής σε MW/Hz (στην κλασική βιβλιογραφία συντελεστής πόλωσης συχνότητας, frequency
bias). Στο απομονωμένο σύστημα ο πρώτος όρος δεν υπάρχει και η \eqref{eq:ace} ανάγεται στο
κριτήριο της προηγούμενης υποενότητας. Για τον λόγο αυτόν ο SO GL ορίζει ως ενιαίο μέγεθος
ελέγχου της αποκατάστασης το \emph{σφάλμα ελέγχου αποκατάστασης συχνότητας} (ΣΕΑΣ·
Frequency Restoration Control Error, FRCE): σε σύστημα πολλών περιοχών ισούται με το ΣΕΠ
της περιοχής, ενώ όταν η περιοχή ελέγχου συμπίπτει με ολόκληρη τη συγχρονισμένη περιοχή
ισούται με την απόκλιση συχνότητας.

Ο συνδυασμός των δύο όρων διακρίνει τη θέση της διαταραχής (Σχήμα~\ref{fig:lfc}):

\begin{itemize}
\item \textbf{Διαταραχή εντός της περιοχής.} Η συχνότητα μειώνεται ($\Delta f < 0$) και οι
      γειτονικές περιοχές παρέχουν ισχύ υποστήριξης, οπότε οι εισαγωγές υπερβαίνουν το
      πρόγραμμα ($\Delta P_{\mathrm{tie}} < 0$). Οι δύο όροι είναι ομόσημοι, το ΣΕΠ λαμβάνει
      μεγάλη απόλυτη τιμή και ενεργοποιείται ΕΑΣ στην περιοχή.
\item \textbf{Διαταραχή εκτός της περιοχής.} Η συχνότητα μειώνεται, αλλά η περιοχή εξάγει
      πρόσθετη ισχύ προς τη γειτονική περιοχή ($\Delta P_{\mathrm{tie}} > 0$). Για κατάλληλη
      τιμή του $K$ οι δύο όροι αλληλοαναιρούνται και το ΣΕΠ παραμένει σχεδόν μηδενικό· η
      περιοχή συμμετέχει μόνο μέσω της ΕΔΣ.
\end{itemize}

Η ιδιότητα αυτή ονομάζεται \emph{μη αλληλεπιδραστικός έλεγχος} (non-interactive control):
κάθε περιοχή αντισταθμίζει μόνο τις δικές της ανισορροπίες, ενώ όλες οι περιοχές
συμμετέχουν μέσω της ΕΔΣ στη συγκράτηση της συχνότητας αμέσως μετά τη διαταραχή. Η
αλληλοαναίρεση είναι πλήρης όταν ο συντελεστής $K$ ισούται με τη χαρακτηριστική
ισχύος-συχνότητας της περιοχής, δηλαδή με το άθροισμα της συνεισφοράς της ΕΔΣ, του
αυτοελέγχου της παραγωγής και της αυτορρύθμισης του φορτίου ανά Hz.

\begin{figure}[!htbp]
\centering
\includegraphics[width=\linewidth]{figures/lfc-areas.svg}
\caption{Εφεδρεία αποκατάστασης συχνότητας σε απομονωμένο και σε διασυνδεδεμένο σύστημα. Στην
πρώτη περίπτωση αρκεί η απόκλιση συχνότητας· στη δεύτερη απαιτείται και ο όρος των
ανταλλαγών ισχύος στις γραμμές διασύνδεσης, ώστε να ενεργοποιείται ΕΑΣ μόνο στην περιοχή
όπου εμφανίστηκε η διαταραχή.}
\label{fig:lfc}
\end{figure}

\subsubsection*{Κατηγορίες εφεδρειών συχνότητας}

\begin{itemize}
\item \textbf{ΕΔΣ} (FCR) --- εφεδρεία διατήρησης συχνότητας. Αυτόματη και τοπική,
      ενεργοποιείται από κοινού σε ολόκληρη τη συγχρονισμένη περιοχή. Συγκρατεί και
      σταθεροποιεί τη συχνότητα.
\item \textbf{αΕΑΣ} (aFRR) --- εφεδρεία αποκατάστασης συχνότητας με αυτόματη ενεργοποίηση.
      Κεντρική, οδηγείται από το ΣΕΑΣ μέσω της ΑΡΠ. Επαναφέρει τη συχνότητα στα 50 Hz και
      αποδεσμεύει την ΕΔΣ.
\item \textbf{χΕΑΣ} (mFRR) --- εφεδρεία αποκατάστασης συχνότητας με χειροκίνητη
      ενεργοποίηση. Ενεργοποιείται με εντολή κατανομής του ΔΣΜ και αποδεσμεύει την αΕΑΣ σε
      παρατεταμένες ανισορροπίες.
\item \textbf{ΕΑ} (RR) --- εφεδρεία αντικατάστασης. Αποκαθιστά το απαιτούμενο επίπεδο των
      ενεργοποιημένων ΕΑΣ, ώστε το σύστημα να είναι έτοιμο για επόμενη διαταραχή.
\end{itemize}

Στην παλαιότερη βιβλιογραφία και σε μέρος της πρακτικής οι τρεις πρώτες κατηγορίες
αναφέρονται ως πρωτεύουσα, δευτερεύουσα και τριτεύουσα ρύθμιση αντιστοίχως. Η ορολογία του
SO GL, με τα ακρωνύμια της επίσημης ελληνικής απόδοσης που χρησιμοποιεί και ο ΑΔΜΗΕ,
κατονομάζει τις εφεδρείες με βάση τη λειτουργία τους (διατήρηση, αποκατάσταση,
αντικατάσταση) και όχι με βάση τη σειρά ενεργοποίησης.

Κάθε επόμενη κατηγορία εφεδρείας ενεργοποιείται βραδύτερα και με χαμηλότερο κόστος από την
προηγούμενη και την αποδεσμεύει, ώστε οι ταχύτερες εφεδρείες να είναι διαθέσιμες για νέα
διαταραχή. Οι χρόνοι πλήρους ενεργοποίησης δίνονται στον Πίνακα~\ref{tab:reserves}.

\begin{table}[!htbp]
\centering
\begin{tabular}{lllc}
\hline
Εφεδρεία & Λειτουργία & Ενεργοποίηση & Χρόνος πλήρους ενεργοποίησης \\
\hline
ΕΔΣ (FCR)   & διατήρηση      & αυτόματη, τοπική   & 30 s \\
αΕΑΣ (aFRR) & αποκατάσταση   & αυτόματη, κεντρική & 5 min \\
χΕΑΣ (mFRR) & αποκατάσταση   & χειροκίνητη        & 12,5 min \\
ΕΑ (RR)     & αντικατάσταση  & χειροκίνητη        & $>$ 15 min \\
\hline
\end{tabular}
\caption{Επίπεδα εφεδρείας συχνότητας στην Ηπειρωτική Ευρώπη, κατά τον SO GL και τα
τυποποιημένα προϊόντα εξισορρόπησης.}
\label{tab:reserves}
\end{table}

\subsubsection*{Παράμετροι συχνότητας για την Ηπειρωτική Ευρώπη}

\begin{itemize}
\item Το \textbf{συμβάν αναφοράς} (reference incident) για τη διαστασιολόγηση της ΕΔΣ είναι
      ανισορροπία 3\,000 MW, είτε ως απώλεια παραγωγής είτε ως απώλεια φορτίου. Η ελάχιστη
      διαθέσιμη ΕΔΣ στη συγχρονισμένη περιοχή είναι ανά πάσα στιγμή 3\,000 MW.
\item Η \textbf{μέγιστη απόκλιση συχνότητας σταθερής κατάστασης} για το συμβάν αναφοράς
      είναι 200 mHz. Από τα δύο μεγέθη προκύπτει ο συντελεστής αναλογίας μεταξύ ΕΔΣ και
      απόκλισης συχνότητας: 3\,000 MW / 0,2 Hz = 15\,000 MW/Hz.
\item Η \textbf{τυπική περιοχή συχνότητας} (standard frequency range) είναι $\pm 50$ mHz· η
      συχνότητα δεν επιτρέπεται να βρίσκεται εκτός αυτής για περισσότερα από
      15\,000 λεπτά ανά έτος. Η μέγιστη στιγμιαία απόκλιση συχνότητας είναι 800 mHz.
\item Απόκλιση συχνότητας σταθερής κατάστασης μεγαλύτερη των 200 mHz κατατάσσει το σύστημα
      σε κατάσταση \emph{έκτακτης ανάγκης}.
\end{itemize}

\subsection{Σύνοψη}

Η λειτουργία του ΣΜΗΕ βασίζεται στον κλειστό βρόχο εποπτείας και ελέγχου που περιγράφηκε
στην αρχή του κεφαλαίου:

\begin{enumerate}
\item Η \textbf{υποδομή εποπτείας} (SCADA, PMU) παρέχει μετρήσεις πλεονάζουσες, με θόρυβο
      και, στην περίπτωση του SCADA, χωρίς χρονικό συγχρονισμό.
\item Η \textbf{εκτίμηση κατάστασης} υπολογίζει από τις μετρήσεις την κατάσταση του δικτύου
      κατά το κριτήριο των σταθμισμένων ελαχίστων τετραγώνων και απορρίπτει τα ανιχνεύσιμα
      εσφαλμένα δεδομένα.
\item Η \textbf{ανάλυση επιχειρησιακής ασφάλειας} αξιολογεί την εκτιμώμενη κατάσταση έναντι
      του κριτηρίου N-1 και κατατάσσει το σύστημα σε μία από τις πέντε καταστάσεις του
      SO GL.
\item Οι \textbf{αποφάσεις λειτουργίας} υλοποιούνται μέσω των αγορών και των μέσων ελέγχου,
      σε χρονικές κλίμακες από τα δευτερόλεπτα της ΕΔΣ έως τους μήνες του
      προγραμματισμού.
\end{enumerate}

Η αξιοπιστία του βρόχου περιορίζεται από την ακρίβεια της εκτιμώμενης κατάστασης: κρίσιμες
μετρήσεις, επιθέσεις συνεπείς με το μοντέλο μετρήσεων και μη παρατηρήσιμα τμήματα του
δικτύου οδηγούν σε εσφαλμένη εκτίμηση χωρίς ένδειξη σφάλματος.

Οι υπηρεσίες που υποστηρίζουν τη λειτουργία του συστήματος ονομάζονται \emph{επικουρικές
υπηρεσίες} και διακρίνονται σε σχετικές και μη σχετικές με τη συχνότητα· εξετάζονται στα
Κεφάλαια 3 και 2 αντιστοίχως. Η λειτουργία των δικτύων διανομής, όπου η παρατηρησιμότητα
είναι σημαντικά χαμηλότερη και η εκτίμηση κατάστασης βασίζεται σε μεγάλο βαθμό σε
ψευδομετρήσεις, εξετάζεται στο Κεφάλαιο 4. Η κυβερνοασφάλεια της υποδομής μέτρησης και
ελέγχου εξετάζεται στο Κεφάλαιο 8.
